\documentclass[12pt,a4paper]{article}
\usepackage[utf8]{inputenc}
\usepackage[russian]{babel}
\usepackage{amsmath}
\usepackage{amsfonts}
\usepackage{listings}
\usepackage{hyperref}

\title{\textbf{ADR 0002: Архитектурный и Научно-Теоретический Журнал Разработки ChainDB}}
\author{Инженерный Совет Проекта ChainDB}
\date{\today}

\begin{document}

\maketitle

\begin{abstract}
Настоящий документ фиксирует фундаментальный научно-теоретический базис проекта \textbf{ChainDB}. В журнале математически обосновано применение Stackful-продолжений в контексте изоморфизма Карри-Ховарда, а также детально описана физика управления памятью рантайма Chez Scheme при реализации низкоуровневой кооперативной многозадачности однопоточного In-Memory хранилища.
\end{abstract}

\section{Статус Документа}
\textbf{Status:} Accepted (Принято к исполнению).

\section{Контекст и Постановка Задачи}
При проектировании высокопроизводительных систем управления базами данных класса In-Memory (таких как Redis), критической уязвимостью является ``проклятие тяжелой команды'' (например, \texttt{KEYS *}). Синхронное выполнение итерационных операций над большими структурами данных намертво блокирует единственный поток ядра, вызывая деградацию всей системы. 

Нам необходимо реализовать механизм кооперативного тайм-шеринга, позволяющий тяжелым операциям добровольно уступать процессорное время быстрым командам (\texttt{GET}/\texttt{SET}) без использования тяжелых потоков ОС (Threads) и без усложнения бизнес-логики явными стейт-машинами.

\section{Научно-Теоретический Багаж Вычислительного Ядра}

\subsection{Изоморфизм Карри-Ховарда и Логика Пирса}
В теории типов и вычислительной топологии ChainDB оператор управления контекстом вычислений \texttt{call/cc} рассматривается не как побочный системный эффект, а как программное воплощение классического закона Пирса:
\begin{equation}
((A \to B) \to A) \to A
\end{equation}

Данная интуиционистская типизация эквивалентна закону снятия двойного отрицания $\neg\neg A \to A$ в классической логике. Наличие данного примитива в рантайме ChainDB позволяет осуществлять нелокальные переходы (``путешествия во времени'') к прошлым состояниям вычислительного стека, восстанавливая инварианты памяти и лениво стримя данные клиентам без накладных расходов на копирование.

\subsection{Физика Памяти Chez Scheme: Spaghetti Stack}
В традиционных Си-подобных архитектурах стек вызовов представляет собой непрерывный монолитный массив памяти. Фиксация контекста в таких системах требует полного копирования байт (\texttt{memcpy}), что приводит к разрушению кэша процессора L1/L2 (\textit{Cache Misses}) и деградации производительности до $O(N)$.

Рантайм Chez Scheme решает эту проблему на уровне архитектуры связанных сегментов стека в куче (\textit{Spaghetti Stack}). При фиксации продолжения через \texttt{call/cc}:
\begin{enumerate}
    \item Текущий активный сегмент стека изолируется и помечается сборщиком мусора (Garbage Collector) как иммутабельный объект в куче ($O(1)$).
    \item Регистр указателя стека процессора (\texttt{RSP/ESP}) перенаправляется на новый пустой фрейм.
    \item При вызове сохраненного продолжения происходит мгновенный сдвиг указателя процессора назад на исходный сегмент в куче, что обеспечивает время переключения контекста $O(1)$.
\end{enumerate}

\section{Реализованные Архитектурные Решения}

\subsection{Компонент: Живые Курсоры (Zero-Copy Iterators)}
Для безопасного итеративного обхода индексов базы данных разработан механизм изолированных фабрик курсоров на базе лексических замыканий. 

Каждый порожденный курсор инкапсулирует внутренние переменные состояния \texttt{resume-cont} (точка сна генератора) и \texttt{return-cont} (точка ожидания пользователя) внутри локального изолированного контейнера \texttt{letrec*}. Это гарантирует полную изоляцию контекстов разных сетевых клиентов и исключает взаимное повреждение памяти (\textit{Memory Corruption}). 

Периметр защищен от бесконечных временных петель (``зомби-генераторов'') за счет принудительного автоматического перевода метки \texttt{resume-cont} в терминальный статус \texttt{'dead} в момент обычного линейного завершения функции пользователя.

\section{Последствия и Риски (Consequences)}
\begin{itemize}
    \item \textbf{Плюсы}: Линейный синтаксис написания тяжелых команд ядра; заморозка контекста работает на любой глубине вложенности функций; переключение задач происходит за субмикросекундное время $O(1)$.
    \item \textbf{Минусы}: Необходимость жесткого контроля точек квантования вычислений (\texttt{yield}). Ошибка разработчика ядра в расстановке пауз может увести главный поток базы в синхронный бесконечный цикл.
\end{itemize}

\end{document}
